%------------------------------------------------------------------------------%

\usepackage{calculo_varias_variaveis-1}

\title{Derivadas de Curvas Paramétricas}

%------------------------------------------------------------------------------%
\begin{document}

\addtitle

%------------------------------------------------------------------------------%
\section{Derivada Primeira}

%------------------------------------------------------------------------------%
\begin{frame}{Fórmula Paramétrica para \(dy/dx\)}
  \pause
  Se $y$ for uma função derivável de $x$ ao longo da curva paramétrica
  \[
    y = y(x)
  \]

  \pause
  Não temos a expressão de $y$ em função de $x$

  \pause
  \vspace{0.2\baselineskip}
  A \emph{regra da cadeia} fornece a relação
  \[
    \frac{dy}{dt} = \frac{dy}{dx} \, \frac{dx}{dt}
  \]

  \pause
  Se \(\frac{dx}{dt}\neq 0\) temos \emph{
  \(
    \frac{dy}{dx} = \frac{\;\sfrac{dy}{dt}\;}{\sfrac{dx}{dt}}
  \)}
\end{frame}

% 03
%------------------------------------------------------------------------------%
\addtocounter{exemplo}{1}
\begin{frame}{Exemplo \theexemplo}
  Encontre a tangente à curva
  \[
    x = \sec(t)
  \]
  \[
    y = \tan(t)
  \]
  \[
    -\frac{\pi}{2} < t < \frac{\pi}{2}
  \]
  no ponto \(\left(\sqrt{2}, 1\right)\), onde \(t=\frac{\pi}{4}\)

  \addgraphic{6}{8,2}{exemplo_calculo_2}
\end{frame}

\begin{frame}{Exemplo \theexemplo}
  Já calculamos
  \[
    \frac{dx}{dt} = \tan(t)\sec(t)
    \qquad
    \frac{dy}{dt} = \sec^2(t)
  \]

  \pause
  Portanto
  \[
    \frac{dy}{dx}
    = \frac{\;\sfrac{dy}{dt}\;}{\sfrac{dx}{dt}}
    \pause
    = \frac{\sec^2(t)}{\tan(t)\sec(t)}
    \pause
    = \frac{\sec(t)}{\tan(t)}
  \]

  \pause
  Avaliando no ponto
  \[
    \frac{dy}{dx} \evalat{t=\sfrac{\pi}{4}}{}
    \pause
    = \frac{\sec\left(\sfrac{\pi}{4}\right)}{\tan\left(\sfrac{\pi}{4}\right)}
    \pause
    = \frac{\sqrt{2}}{1}
    \pause
    = \sqrt{2}
  \]
\end{frame}

\begin{frame}{Exemplo \theexemplo}
  Equação da reta
  \begin{align*}
    \onslide<+->{y - y_0 & = a (x - x_0)                                  }
    \onslide<+->{          = \frac{dy}{dx}\evalat{t_0}{} (x - x_0) \\[3mm]}
    \onslide<+->{y - 1   & = \sqrt{2} \left(x - \sqrt{2}\right)           }
    \onslide<+->{          = \sqrt{2} x - 2                        \\[3mm]}
    \onslide<+->{y       & = \sqrt{2} x - 2 + 1                           }
    \onslide<+->{          = \sqrt{2} x - 1                               }
  \end{align*}

  \onslide<+->{
    Equação da reta tangente
    \[
      y = \sqrt{2} x - 1
    \]}
\end{frame}

% 04
%------------------------------------------------------------------------------%
\addtocounter{exemplo}{1}
\begin{frame}{Exemplo \theexemplo}
	Encontre uma equação para a reta tangente\\
	ao astróide
	\[
	x = \cos^3(t) \qquad
	y = \sin^3(t) \qquad
	0 \leq t \leq 2\pi
	\]
	no ponto \(t=\frac{\pi}{6}\)
	\addgraphic{6}{9.5,2}{exemplo_calculo_astroide}
\end{frame}

\begin{frame}{Exemplo \theexemplo}
	\[
		\frac{dx}{dt}
		\pause
		= \frac{d}{dt}\cos^3(t)
		\pause
		= 3\cos^2(t)\frac{d}{dt}\cos(t)
		\pause
		= -3\cos^2(t)\sin(t)
	\]

	\pause
	\[
		\frac{dy}{dt}
		\pause
		= \frac{d}{dt}\sin^3(t)
		\pause
		= 3\sin^2(t)\frac{d}{dt}\sin(t)
		\pause
		= 3\sin^2(t)\cos(t)
	\]

	\pause
	Portanto
	\[
		\frac{dy}{dx}
		\pause
		= \frac{\;\sfrac{dy}{dt}\;}{\sfrac{dx}{dt}}
		\pause
		= \frac{3\sin^2(t)\cos(t)}{-3\cos^2(t)\sin(t)}
		\pause
		= -\frac{\sin(t)}{\cos(t)}
		\pause
		= -\tan(t)
	\]

	\pause
	Avaliando no ponto
	\[
		\frac{dy}{dx} \evalat{t=\sfrac{\pi}{6}}{}
		\pause
		= -\tan\left(\frac{\pi}{6}\right)
		\pause
		= -\frac{\sqrt{3}}{3}
		\pause
		= \frac{-1}{\sqrt{3}}
		= \frac{-1}{3^{\sfrac{1}{2}}}
	\]
\end{frame}

\begin{frame}{Exemplo \theexemplo}
	Equação da reta
	\[
		y - y_0
		= a (x - x_0)
		\pause
		= \frac{dy}{dx}\evalat{t_0}{} (x - x_0)
	\]

	\pause
	\[
		x_0
		\pause
		= x\left(\frac{\pi}{6}\right)
		\pause
		= \cos^3\left(\frac{\pi}{6}\right)
		\pause
		= \left(\frac{\sqrt{3}}{2}\right)^3
		\pause
		= \frac{3^{\sfrac{3}{2}}}{2^3}
	\]

	\pause
	\[
		y_0
		\pause
		= y\left(\frac{\pi}{6}\right)
		\pause
		= \sin^3\left(\frac{\pi}{6}\right)
		\pause
		= \left(\frac{1}{2}\right)^3
		\pause
		= \frac{1}{2^3}
	\]
\end{frame}

\begin{frame}{Exemplo \theexemplo}
	Equação da reta
	\begin{align*}
		\onslide<+->{y - y_0 & = \frac{dy}{dx}\evalat{t_0}{} (x - x_0) \\[3mm]}
		\onslide<+->{y - \frac{1}{2^3}
			         & = \frac{-1}{3^{\sfrac{1}{2}}} \left(x - \frac{3^{\sfrac{3}{2}}}{2^3}\right) }
		\onslide<+->{  = \frac{-x}{3^{\sfrac{1}{2}}}
			           - \frac{-1\;}{3^{\sfrac{1}{2}}} \frac{3^{\sfrac{3}{2}}}{2^3} \\[3mm]}
		\onslide<+->{y & = -\frac{x}{\sqrt{3}} + \frac{3}{2^{3}} + \frac{1}{2^3}       }
		\onslide<+->{    = -\frac{x}{\sqrt{3}} + \frac{4}{8} }
		\onslide<+->{    = -\frac{x}{\sqrt{3}} + \frac{1}{2} }
	\end{align*}

	\onslide<+->{
		Equação da reta tangente
		\[
		y = -\frac{x}{\sqrt{3}} + \frac{1}{2}
		\]}
\end{frame}

%------------------------------------------------------------------------------%
\section{Derivada Segunda}

%------------------------------------------------------------------------------%
\begin{frame}{Fórmula Paramétrica para \(d^2y/dx^2\)}
  \begin{columns}
    \column{0.5\linewidth}
      Sabemos que
      \[
        \frac{dy}{dx} = \frac{\;\dfrac{dy}{dt}\;}{\dfrac{dx}{dt}}
      \]

      \pause
      Substituindo $y$ por $y'$ temos
      \[
        \frac{d\emph{y'}}{dx} = \frac{\;\dfrac{d\emph{y'}}{dt}\;}{\dfrac{dx}{dt}}
      \]
    \column{0.5\linewidth}
      \pause
      \emph{
      \[
        \frac{d^2y}{dx^2}
        \pause
        = \frac{d y'}{dx}
        \pause
        = \frac{\;\dfrac{dy'}{dt}\;}{\dfrac{dx}{dt}}
      \]}
  \end{columns}
\end{frame}

% 05
%------------------------------------------------------------------------------%
\addtocounter{exemplo}{1}
\begin{frame}{Exemplo \theexemplo}
  Encontre \(\sfrac{d^2y}{dx^2}\) como função de $t$ se
  \[
    x = t - t^2 \qquad
    y = t - t^3
  \]
\end{frame}

\begin{frame}{Exemplo \theexemplo}
  Vamos usar a fórmula
  \(
    \frac{d^2y}{dx^2}
    = \frac{\;\sfrac{dy'}{dt}\;}{\sfrac{dx}{dt}}
  \)

  \pause
  Calculando \(\sfrac{dx}{dt}\) e \(\sfrac{dy}{dt}\)
  \pause
  \[
    \frac{dx}{dt}
    \pause
    = \frac{d}{dt}\left(t - t^2\right)
    \pause
    = 1 - 2t
  \]
  \pause
  \[
    \frac{dy}{dt}
    = \frac{d}{dt}\left(t - t^3\right)
    \pause
    = 1 - 3t^2
  \]

  \pause
  Calculando \(\sfrac{dy}{dx}\)
  \[
    \pause
    y'
    = \frac{dy}{dx}
    \pause
    = \frac{\sfrac{dy}{dt}}{\sfrac{dx}{dt}}
    \pause
    = \frac{1 - 3t^2}{1 - 2t}
  \]
\end{frame}

\begin{frame}{Exemplo \theexemplo}
  \begin{align*}
    \onslide<+->{\frac{dy'}{dt}
    & = \frac{d}{dt}\left(\frac{1 - 3t^2}{1 - 2t}\right) \\[2mm]}
    \onslide<+->{& = \frac{\frac{d}{dt}\left(1 - 3t^2\right)\left(1 - 2t\right) - \left(1 - 3t^2\right)\frac{d}{dt}\left(1 - 2t\right)}
             {\left(1 - 2t\right)^2} \\[2mm]}
    \onslide<+->{& = \frac{\left(-6t\right)\left(1 - 2t\right) - \left(1 - 3t^2\right)\left(-2\right)}
             {\left(1 - 2t\right)^2} \\[2mm]}
    \onslide<+->{& = \frac{\left(-6t +12t^2\right) - \left(-2 +6t^2\right)}
             {\left(1 - 2t\right)^2} \\[2mm]}
    \onslide<+->{& = \frac{2 -6t +6t^2}
             {\left(1 - 2t\right)^2}}
  \end{align*}
\end{frame}

\begin{frame}{Exemplo \theexemplo}
  \[
    \onslide<+->{\frac{d^2y}{dx^2}}
    \onslide<+->{= \frac{\;\dfrac{dy'}{dt}\;}{\dfrac{dx}{dt}}}
    \onslide<+->{= \frac{\;\dfrac{2 -6t +6t^2}{\left(1 - 2t\right)^2}\;}{1 - 2t}}
    \onslide<+->{= \frac{2 -6t +6t^2}{\left(1 - 2t\right)^3}}
  \]
\end{frame}

% 06
%------------------------------------------------------------------------------%
\addtocounter{exemplo}{1}
\begin{frame}{Exemplo \theexemplo}
	Calcule \(\sfrac{dy}{dx}\), em $t = 2$,
	para a curva paramétrica definida implicitamente
	\[
		x^3 + 2t^2 = 9
		\qquad
		2y^3 - 3t^2 = 4
	\]

	\pause
	Nesse caso podemos isolar $x$ e $y$
	\[
	  x = \sqrt[3]{9-2t^2}
	  \qquad
	  y = \sqrt[3]{\frac{4+3t^2}{2}}
	\]
	\pause
	Se não for possível?
\end{frame}

\begin{frame}{Exemplo \theexemplo}
\begin{columns}
\column{0.5\linewidth}
	\begin{align*}
		\onslide<+->{x^3 + 2t^2 & = 9 \\[3mm]}
		\onslide<+->{x^3 & = 9 - 2t^2  \\[3mm]}
		\onslide<+->{\frac{d}{dt}\left(x^3\right) & = \frac{d}{dt}\left(9 - 2t^2\right) \\[3mm]}
		\onslide<+->{3x^2\frac{dx}{dt} & = 0 - 2\times 2t \\[3mm]}
		\onslide<+->{\frac{dx}{dt} & = \frac{-4t}{3x^2}}
	\end{align*}
\column{0.5\linewidth}
	\begin{align*}
		\onslide<+->{2y^3 - 3t^2 & = 4 \\[3mm]}
		\onslide<+->{2y^3 & = 4 + 3t^2 \\[3mm]}
		\onslide<+->{\frac{d}{dt}\left(2y^3\right) & = \frac{d}{dt}\left(4 + 3t^2\right) \\[3mm]}
		\onslide<+->{2\times 3y^2\frac{dy}{dt} & = 0 + 3 \times 2 t \\[3mm]}
		\onslide<+->{\frac{dy}{dt} & = \frac{6t}{6y^2}}
		\onslide<+->{= \frac{t}{y^2}}
	\end{align*}
	\end{columns}
\end{frame}

\begin{frame}{Exemplo \theexemplo}
\[
	\frac{dy}{dx}
	\pause
	= \frac{\;\dfrac{dy}{dt}\;}{\dfrac{dx}{dt}}
	\pause
	= \frac{\dfrac{t}{y^2}}{\;\dfrac{\;-4t\;}{3x^2}\;}
	\pause
	= \frac{t}{y^2} \; \frac{3x^2}{\;-4t\;}
	\pause
	= -\frac{3x^2}{4y^2}
\]
\end{frame}

\begin{frame}{Exemplo \theexemplo}
Avaliando $x_0$ e $y_0$ no ponto $t_0=2$
\begin{columns}
	\column{0.5\linewidth}
	\begin{align*}
	\onslide<+->{x_0^3 + 2t_0^2     & = 9 \\}
	\onslide<+->{x_0^3 + 2\times2^2 & = 9 \\}
	\onslide<+->{x_0^3 + 8          & = 9 \\}
	\onslide<+->{x_0^3              & = 1 \\}
	\onslide<+->{x_0                & = 1}
\end{align*}
\column{0.5\linewidth}
\begin{align*}
	\onslide<+->{2y_0^3 - 3t_0^2      & = 4  \\}
	\onslide<+->{2y_0^3 - 3\times 2^2 & = 4  \\}
	\onslide<+->{2y_0^3 - 12          & = 4  \\}
	\onslide<+->{2y_0^3               & = 16 \\}
	\onslide<+->{ y_0^3               & = 8  \\}
	\onslide<+->{ y_0                 & = 2}
\end{align*}
\end{columns}
\end{frame}

\begin{frame}{Exemplo \theexemplo}
\[
	\frac{dy}{dx} \evalat{t=2}{}
	\pause
	= -\frac{3x^2}{4y^2}\evalat{t=2}{}
	\pause
	= -\frac{3x_0^2}{4y_0^2}
	\pause
	= -\frac{3\times 1^2}{4\times 2^2}
	\pause
	= -\frac{3}{16}
\]
\end{frame}

%------------------------------------------------------------------------------%
\section{Lista Mínima}

%------------------------------------------------------------------------------%
\begin{frame}{Lista Mínima}
  Cálculo Vol. 2 do Thomas 12\textsuperscript{a} ed. -- Seção 11.2
  \begin{enumerate}
    \item Estudar todo o texto da seção
    \item Resolver os exercícios: 2, 8, 12, 16, 18, 20
  \end{enumerate}

  \vfill
  Atenção: A prova é baseada no livro, não nas apresentações
\end{frame}

%------------------------------------------------------------------------------%
\end{document}
%------------------------------------------------------------------------------%
